%=================================================================
\section{Data Availability Details}
\label{app:data}

% src: data/backups/gex_llm/condensed/ (row counts read from the files, as of 2026-11-10; 06_condense.py: detections 73,846 in all, of which 36,706 are answers of gpt-6-luna and o4-mini (34,784 + 1,922); windows 38,592; series 16; runs 179 in all, of which 46 are Batch jobs, the synchronous pilot and backfilled runs (43 gpt-6-luna + 3 o4-mini)); series_index.csv; runs.csv.gz and run_records.jsonl.gz (06_condense.py, since 2026-11-09, D-17 revision); scripts/validation/electronics/README.md (storage section); 02_decisions_log.md D-14, D-17, D-31 (raw chains outside the deposit because of their size)
The repository holds the condensed material and the results. The condensed files are in \path{data/backups/gex_llm/condensed/}. The files and the full tables are shared with a companion paper that repeats seven of the arms with open-weights models on local hardware \cite{regan2026openweights}: the deposit also holds those answers and runs, which this article does not use. The counts below are those of the hosted and backfilled arms of this article, as of 2026-11-10; they do not include the answers and runs of the companion paper, and the totals of the files grow when its runs are stored.
\begin{itemize}[leftmargin=*,labelsep=5.5mm]
\item \path{detections.csv.gz}: 36{,}706 rows for this article (the file also holds the answers of the companion paper), one per stored model answer, with the verdict, the reported metrics and the parse status but without the reasoning or the raw text. The key is \ident{arm}, \ident{series_key}, \ident{symbol}, \ident{window_end}, \ident{model}, \ident{prompt_version}.
\item \path{windows.csv.gz}: 38{,}592 rows, one per 30-day window, with the labels, the sha256 of the input as sent and the verdict of the deterministic rule but without the sequences.
\item \path{series_index.csv}: 16 rows, one per series and symbol, with the row count, the first and last date, the source table and the hash of the builder.
\item \path{runs.csv.gz}: 46 rows for this article, one per Batch job, pilot or backfilled run (and further rows for the local runs of the companion paper), with the batch and file identifiers, the sha256 of the input and output files, the token counts and the cost.
\item \path{run_records.jsonl.gz}: the manifest and the batch record of every run folder (\path{manifest.json} and \path{batch_<k>.json}), one line per file.
\end{itemize}
Every detection rate and agreement figure can be recomputed from the first two files (join on \ident{series_key}, \ident{symbol} and \ident{window_end}); the other numbers of the article are in the files of \path{docs/papers/electronics/results/}, and the daily values behind them are in the full tables. The provenance of every Batch job is in the last two files, which are written from the per-run files (\path{tables/runs/}, and \path{manifest.json} and \path{batch_<k>.json} in each run folder). The per-run files, the raw request and answer files (\path{input_<k>.jsonl}, \path{output_<k>.jsonl}) and the full tables are kept on the authors' machine and in the database; they are not part of the repository, and the full tables are deposited at Zenodo.

The full tables are loaded into four tables and one view of the schema \ident{backups} of the research database: \ident{gex_llm_regime_series} (the number shown to the model, one row per series, symbol and trading date), \ident{gex_llm_regime_windows} (the 30-day input as sent, its sha256 and the verdict of the rule), \ident{gex_llm_regime_runs} (one row per Batch job, pilot or backfilled run, and per local run of the companion paper) and \ident{gex_llm_regime_detections} (one row per model answer), with the view \ident{gex_llm_regime_rates_summary} (yearly rate of the model next to the rate of the rule, by arm). The parquet copies of these four tables are deposited at Zenodo (Data Availability Statement). The raw option chains are not part of the deposit because of their size; it holds the daily aggregates that were shown to the models.
